\documentclass[10pt,a4paper]{article}
\usepackage[utf8]{inputenc}
\usepackage[T1]{fontenc}
\usepackage{lmodern}
\usepackage[a4paper,margin=15mm]{geometry}
\usepackage{amsmath,amssymb,booktabs,graphicx}
\usepackage{enumitem}
\usepackage{xcolor}
\usepackage{hyperref}
\usepackage{fancyhdr}
\usepackage[round,authoryear]{natbib}
\usepackage{titlesec}
\usepackage{array}
\usepackage{url}
\usepackage{setspace}
\usepackage{placeins}
\usepackage{float}
\hypersetup{colorlinks=true,linkcolor=blue,citecolor=blue,urlcolor=blue}
\bibpunct{(}{)}{; }{a}{}{, }
\renewcommand{\familydefault}{\rmdefault}
\urlstyle{same}
\setlength{\parindent}{1em}
\setlength{\parskip}{0pt}
\setlength{\textfloatsep}{2pt plus 1pt minus 1pt}
\setlength{\dbltextfloatsep}{2pt plus 1pt minus 1pt}
\setlength{\intextsep}{2pt plus 1pt minus 1pt}
\setlength{\floatsep}{3pt plus 1pt minus 1pt}
\setstretch{0.97}
\AtBeginDocument{\fontsize{8}{8.7}\selectfont}
\titleformat{\section}{\normalfont\bfseries\normalsize}{\thesection}{0.5em}{}
\titlespacing*{\section}{0pt}{5pt}{2pt}
\titleformat{\subsection}{\normalfont\bfseries\small}{\thesubsection}{0.5em}{}
\titlespacing*{\subsection}{0pt}{4pt}{1pt}
\pagestyle{fancy}
\fancyhf{}
\fancyhead[L]{AI-assisted CT and MRI triage}
\fancyhead[R]{Article paper}
\fancyfoot[C]{\thepage}
\newcommand{\keywords}[1]{\par\noindent\textbf{Keywords:} #1}
\newcommand{\studyresult}[1]{\textit{#1}}

\title{\textbf{A Computer Vision Framework Combining Machine Learning and Deep Learning for Multimodal Triage of Stroke, Hemorrhage, and Mass Effect on Brain CT and MRI}}
\author{Vivian Moraa Nyamari\\
\small Strathmore University -- @iLabAfrica, Institute of Mathematical Sciences\\
\small \href{mailto:vnyamari@strathmore.edu}{vnyamari@strathmore.edu}\\
\small \href{mailto:viviannyamoraa@gmail.com}{viviannyamoraa@gmail.com}}
\date{}

\begin{document}
\twocolumn[
\maketitle
\begin{abstract}
Acute stroke, intracranial hemorrhage, and traumatic brain injury require prompt assessment because delays in recognizing critical findings can affect treatment decisions and patient outcomes. However, timely interpretation of brain imaging may be difficult in resource-constrained settings where radiology expertise, specialist coverage, and diagnostic infrastructure are limited. This article presents the design of a computer-vision triage framework for brain computed tomography (CT) and magnetic resonance imaging (MRI), bringing data science and analytics together with clinical medicine and surgery. Its objective is to support the prioritization of stroke, hemorrhage, mass effect, and apparently normal studies without replacing professional diagnosis. The workflow accepts CT images, MRI images, or both CT and MRI images within a single study, while retaining the modality assigned to each image. It applies documented image validation, modality-aware preprocessing, and quantitative feature extraction, and defines a comparative modeling plan for feature-based machine learning and convolutional-neural-network deep learning. The application is designed to combine image-derived screening impressions with clinician-entered presenting history and examination information, then present structured findings, confidence, priority, and differential considerations for qualified review. Evaluation is specified using patient-level data partitions and clinically relevant measures of discrimination, sensitivity, specificity, F1-score, area under the receiver-operating-characteristic curve, calibration, subgroup performance, and robustness across modalities, sites, scanners, and acquisition conditions. The framework also incorporates explainability, privacy safeguards, data governance, human oversight, and implementation considerations for African and other settings with constrained resources. The current application is a research prototype, and no clinical performance or outcome benefit is claimed. Such claims require representative, appropriately annotated datasets, external validation, and prospective workflow evaluation. This work therefore contributes a reproducible design and validation pathway for responsible AI-assisted neuroimaging triage.
\end{abstract}
\keywords{brain imaging; computed tomography; magnetic resonance imaging; stroke; hemorrhage; mass effect; deep learning; explainable AI; triage; clinical decision support}
\vspace{3mm}
]

\section{Introduction}
Stroke is a leading cause of death and disability worldwide, and the clinical value of treatment is strongly time dependent \citep{who2022,feigin2022,campbell2020}. Emergency neuroimaging is therefore not merely a diagnostic activity; it is also a prioritization mechanism. Non-contrast CT is widely available and fast, while MRI supplies complementary information about tissue injury, diffusion restriction, perfusion, and lesion age \citep{powers2019,wintermark2013,chalela2007}. Intracranial hemorrhage and mass effect represent particularly important urgent findings because they may require immediate neurosurgical, neurocritical-care, or blood-pressure management decisions \citep{greenberg2022,hemphill2001}.

Radiology workloads, unequal access to subspecialty expertise, and the increasing volume of imaging create a practical need for computer-assisted triage. A triage model does not need to make the final diagnosis to be useful: it may identify studies that should be reviewed sooner, summarize image-derived evidence, and make the reasoning pathway auditable. However, clinical deployment requires careful separation of screening from diagnosis, explicit uncertainty, robust external validation, and human oversight \citep{topol2019,kelly2019,holzinger2019}.

This article develops a combined CT--MRI research and software framework. Its contributions are:
\begin{enumerate}[leftmargin=*,nosep]
\item A single workflow that retains modality identity for every uploaded image.
\item A comparison between feature-based machine learning and image-based deep learning.
\item Clinician-oriented outputs for stroke, hemorrhage, mass effect, and normal categories.
\item A reproducible evaluation and governance plan.
\end{enumerate}

\section{Clinical and Research Background: Literature Review}
\subsection{Clinical context and imaging roles}
Ischemic stroke results from interrupted cerebral perfusion, while hemorrhagic stroke involves bleeding within or around the brain. Traumatic brain injury can cause hemorrhage, edema, contusion, and secondary mass effect. Mass effect describes displacement or compression of intracranial structures and may arise from several underlying conditions. These findings may coexist; therefore, a four-class label is a useful prototype simplification, not a complete representation of clinical diagnosis.

Non-contrast CT is fast and widely available in emergency care and is central to identifying acute intracranial blood. MRI, including diffusion-weighted sequences, can detect early ischemic tissue injury and provide complementary characterization, but acquisition time, scanner access, and patient stability affect its use \citep{chalela2007,warach1992,kidwell2004,wintermark2013}. The clinical implication is that CT and MRI should be treated as complementary, modality-specific inputs. A model must not infer that the absence of a finding on one modality guarantees its absence on the other.

\subsection{Literature review approach}
This narrative review synthesizes clinical imaging guidance, representative neuroimaging-AI studies, foundational computer-vision methods, and work on explainability and clinical implementation. The cited works are used to identify methods and research gaps; this is not a registered systematic review or a quantitative meta-analysis. Study selection for a final dissertation should be documented with search dates, databases, exact query strings, inclusion criteria, and screening counts. Google Scholar and ResearchGate are useful discovery and access services, but study metadata and claims should be verified against the publisher record or bibliographic database before submission.

\subsection{Representative prior work}
Earlier research demonstrates that deep-learning systems can detect selected urgent findings in head CT, while broader medical-imaging reviews describe the strengths and limitations of convolutional models \citep{titano2018,chilamkurthy2018,kuo2019,arbabshirani2018,litjens2017,shen2017}. These studies establish feasibility for defined tasks; they do not, by themselves, validate the present four-label, CT--MRI, mixed-study prototype or demonstrate benefit in the intended local setting. Table~\ref{tab:prior-work} summarizes the relevance and remaining gap without treating results from different datasets as directly comparable.

\begin{table*}[t]
\centering
\small
\caption{Representative prior work informing the proposed study. The table summarizes study focus, not comparative performance; datasets and endpoints differ across publications.}
\label{tab:prior-work}
\begin{tabular}{p{0.15\textwidth}p{0.20\textwidth}p{0.30\textwidth}p{0.28\textwidth}}
\toprule
Study & Imaging/task focus & Contribution relevant to this thesis & Limitation or remaining question addressed here \\
\midrule
\citep{chalela2007} & Emergency stroke imaging; CT and MRI & Compared the role of CT and MRI in acute stroke assessment. & Clinical comparison does not supply a validated multimodal AI triage pipeline. \\
\citep{chilamkurthy2018} & Critical findings on head CT & Demonstrated deep-learning analysis for several urgent CT findings. & Focus is CT; local transferability and the present label set require separate evaluation. \\
\citep{titano2018} & Cranial imaging surveillance & Investigated automated prioritization of acute neurological events. & Workflow and dataset context differ; prospective benefit cannot be assumed for this prototype. \\
\citep{kuo2019} & Intracranial hemorrhage on head CT & Evaluated deep-learning support for a focused hemorrhage-detection task. & Single-finding CT performance does not establish four-class, cross-modality performance. \\
\citep{arbabshirani2018} & Intracranial hemorrhage on brain CT & Examined machine learning for automated CT identification. & External validation, class definitions, and workflow outcomes remain context dependent. \\
\citep{litjens2017,shen2017} & Reviews of deep learning in medical imaging & Summarized architectures, applications, and recurring data challenges. & Reviews motivate a task-specific protocol but do not validate the current software or its outputs. \\
\bottomrule
\end{tabular}
\end{table*}
\FloatBarrier

\subsection{Evidence synthesis and research gap}
Figure~\ref{fig:evidence-gap} presents a qualitative synthesis of the literature. It is a conceptual evidence map, not a numerical graph of publication counts or model accuracy. The map shows how the present thesis connects established, task-specific imaging-AI work to an explicit multimodal prototype and a validation plan.

\begin{figure}[H]
\vspace*{2pt}
\centering
\begin{tabular}{c}
\fbox{\begin{minipage}{0.92\columnwidth}\centering\scriptsize\textbf{Prior evidence}\par CT critical-finding AI; CT/MRI clinical comparisons; imaging-AI reviews.\end{minipage}}\\[-2pt]
$\Downarrow$\\[-1pt]
\fbox{\begin{minipage}{0.92\columnwidth}\centering\scriptsize\textbf{Research gap}\par Evidence is split by modality/task; local robustness and clinician-facing outputs need study.\end{minipage}}\\[-2pt]
$\Downarrow$\\[-1pt]
\fbox{\begin{minipage}{0.92\columnwidth}\centering\scriptsize\textbf{This thesis}\par CT, MRI, or both; compare feature ML and CNN; evaluate transparently.\end{minipage}}
\end{tabular}
\caption{Qualitative literature-to-thesis evidence map. Arrows indicate research motivation, not proven causal or performance improvement.}
\label{fig:evidence-gap}
\end{figure}
\FloatBarrier

\subsection{Why this thesis is important}
The thesis is important for three connected reasons. Clinically, it focuses on prioritizing time-sensitive brain findings while keeping a professional responsible for final interpretation. Methodologically, it compares interpretable engineered features with image-based deep-learning approaches and requires patient-level and external evaluation rather than relying on accuracy alone. From a data-science and health-systems perspective, it makes modality identity, data quality, uncertainty, governance, and resource constraints explicit design variables. The study's contribution is a reproducible evaluation pathway and a testable prototype, not a claim that the current system is ready for diagnosis or has already improved outcomes.

\section{Problem Statement and Objectives}
\subsection{Problem statement}
Many existing studies address a single modality or a single finding. A practical emergency tool must handle heterogeneous studies, preserve CT/MRI provenance, and communicate its output in language that supports review. The problem addressed here is the design of an auditable, multimodal, multi-class triage system for stroke, hemorrhage, mass effect, and normal studies.

\subsection{General objective}
To develop and evaluate a multimodal clinical decision-support prototype that analyzes brain CT and MRI images and assigns a clinically interpretable triage category for stroke, hemorrhage, mass effect, and normal findings.

\subsection{Specific objectives}
\begin{enumerate}[leftmargin=*,nosep]
\item To design a data pipeline that preserves CT and MRI modality identity throughout ingestion, preprocessing, modeling, and reporting.
\item To compare feature-based machine-learning models with image-based deep-learning models.
\item To evaluate sensitivity, specificity, precision, F1-score, AUROC, calibration, and workload implications.
\item To develop clinician-oriented outputs using feature attribution and image-saliency explanations.
\item To assess model stability across sites, scanners, protocols, and image-quality conditions.
\item To define reproducible governance, validation, monitoring, and human-oversight requirements for safe use.
\end{enumerate}

\subsection{Research questions}
The research questions are:
\begin{enumerate}[leftmargin=*,nosep]
\item How does performance differ between CT and MRI?
\item How do radiomic and deep representations compare?
\item Which findings are most reliably prioritized?
\item How stable are predictions under site and scanner variation?
\item Can a concise result card improve the speed and clarity of clinical review without encouraging automation bias?
\end{enumerate}

\section{Materials and Methods}
\subsection{Study design}
This is a retrospective, multi-stage development study followed by external and prospective evaluation. The development phase is intended to use de-identified images organized into training, validation, and held-out test partitions. Patient-level splitting is mandatory to avoid leakage between slices or studies from the same patient. A proposed initial partition is 70\% training, 15\% validation, and 15\% test, subject to feasibility, site representation, and sufficient class counts. The final test set must remain inaccessible until model selection and threshold specification are complete. If the available dataset is too small for a stable three-way split, grouped cross-validation should be used during development and an independent external cohort reserved for final evaluation.

\subsection{Data sources and labeling}
Candidate data may be obtained from institutionally approved archives, public research datasets, and controlled-access repositories. No dataset size or class distribution is asserted in this article because a curated, labeled study cohort has not yet been documented in the workspace. Each case should have a de-identified patient and study identifier, modality and sequence, acquisition metadata where available, label provenance, and quality-control status. Labels should be defined in a written annotation manual, assigned by qualified readers, and adjudicated for disagreement. Weak labels derived from reports must be identified separately from expert annotations. Data inclusion and exclusion counts, missingness, duplicate handling, and reasons for exclusion must be reported in the final results.

The four prototype labels are:
\begin{itemize}[leftmargin=*,nosep]
\item \emph{stroke}
\item \emph{hemorrhage}
\item \emph{mass effect}
\item \emph{normal}
\end{itemize}
Multi-label cases should be retained in the source registry even if the initial classifier uses a mutually exclusive output. A future multilabel model may represent coexisting abnormalities more faithfully.

\subsection{Preprocessing}
The preprocessing pipeline standardizes orientation and image size, converts supported formats to a controlled representation, normalizes intensities, and records every transformation. CT preprocessing may include a documented windowing policy for brain and blood-sensitive views. MRI preprocessing may include bias-field correction, resampling, intensity standardization, and sequence harmonization when metadata are available. No transformation may use information from the held-out test set.

Images with severe motion, missing slices, corrupted pixels, or incompatible metadata are flagged rather than silently discarded. The software records the original path, transformed path, modality, dimensions, and quality-control decision. Data augmentation is applied only to training data and may include small rotations, translations, intensity changes, and noise consistent with clinical acquisition.

\subsection{Feature engineering}
The feature-based branch computes intensity statistics, entropy, texture descriptors, gradient or edge density, local contrast, and shape or region summaries when segmentation is available. Let $x$ denote the feature vector and $y$ the clinical class. A standardized feature matrix is formed as
\begin{equation}
z_j = \frac{x_j-\mu_j}{\sigma_j+\epsilon},
\end{equation}
where the mean and standard deviation are fitted on the training partition only. Feature engineering is evaluated for stability because a feature that changes substantially across scanners may create a misleading shortcut.

\subsection{Machine-learning models}
The classical branch compares logistic regression, support-vector machines, random forests, and gradient-boosted trees. Hyperparameters are selected using nested cross-validation on the training data. Class weighting or balanced sampling is considered when prevalence is unequal. The final model is serialized together with its feature schema, preprocessing parameters, training data version, and software environment.

\subsection{Deep-learning models}
The image branch evaluates transfer-learning backbones such as ResNet, DenseNet, and EfficientNet. A baseline model uses a frozen feature extractor followed by a task-specific head; a second configuration fine-tunes selected upper layers. The loss function may use weighted cross-entropy or focal loss when rare urgent findings are underrepresented. Early stopping is based on validation loss and clinically relevant metrics rather than training accuracy alone.

\subsection{Proposed model structure and comparison}
The planned comparison contains two alternatives, not an assumed fused clinical model. The feature-based branch receives the engineered feature vector and uses a standardized multiclass classifier. A transparent baseline is multinomial logistic regression trained with softmax cross-entropy; random forest and support-vector models may be added only if their dependencies and evaluation are reproducible. For input feature vector $x\in\mathbb{R}^{p}$, the class probability for class $k$ is
\begin{equation}
P(y=k\mid x)=\frac{\exp(w_k^{\mathsf T}x+b_k)}{\sum_{j=1}^{K}\exp(w_j^{\mathsf T}x+b_j)},\qquad K=4.
\end{equation}
The image-based branch receives a preprocessed image tensor and uses a 2-D convolutional backbone initialized from public pretrained weights when licensing and domain suitability permit. The backbone output is globally pooled and passed to a task-specific dense layer, dropout regularization, and a four-class softmax head. CT and MRI are routed with modality metadata and evaluated separately; a late-fusion model that combines examinations is exploratory and must not be reported as validated unless separately trained and tested on paired patient studies.

Figure~\ref{fig:model-structure} summarizes the planned comparative architecture. It shows data flow and candidate model components; it does not mean the CNN branch is already implemented or trained.

\begin{figure}[!t]
\vspace*{-8pt}
\centering
\begin{tabular}{c}
\fbox{\begin{minipage}{0.92\columnwidth}\centering\scriptsize\textbf{1. Study ingestion}\quad CT, MRI, or both; preserve study and modality labels.\end{minipage}}\\[-2pt]
$\Downarrow$\\[-1pt]
\fbox{\begin{minipage}{0.92\columnwidth}\centering\scriptsize\textbf{2. Validation and preprocessing}\quad Check files and quality; normalize by modality and resize.\end{minipage}}\\[-2pt]
$\Downarrow$\\[-1pt]
\fbox{\begin{minipage}{0.92\columnwidth}\centering\scriptsize\textbf{3. Comparative models}\quad Feature ML: vector $\rightarrow$ scaling $\rightarrow$ softmax. Image DL: image $\rightarrow$ CNN $\rightarrow$ pooled features $\rightarrow$ class head.\end{minipage}}\\[-2pt]
$\Downarrow$\\[-1pt]
\fbox{\begin{minipage}{0.92\columnwidth}\centering\scriptsize\textbf{4. Evaluation and review}\quad Metrics, explanations, clinician-facing result, professional review.\end{minipage}}
\end{tabular}
\caption{Proposed comparative model architecture. The CNN branch is planned research work; the current running image-triage path uses heuristic feature scores, as described in the prototype-status results.}
\label{fig:model-structure}
\end{figure}
\FloatBarrier

Table~\ref{tab:model-branches} clarifies the inputs, structure, outputs, and current development status of each path.
\begin{table*}[t]
\centering
\small
\caption{Proposed model branches and current implementation status.}
\label{tab:model-branches}
\begin{tabular}{p{0.15\textwidth}p{0.20\textwidth}p{0.29\textwidth}p{0.29\textwidth}}
\toprule
Branch & Input & Model structure & Status and required evidence \\
\midrule
Prototype image triage & One normalized grayscale image plus CT/MRI label & Six summary image features, heuristic class scores, fixed priority mapping & Implemented as a software prototype; not learned from a labeled cohort and not clinically validated. \\
Feature-based ML & Eight engineered image features from dataset loader & Training-only scaling, multinomial softmax/logistic baseline; optional comparator models & Training code exists; correct split use, non-placeholder held-out metrics, calibration, and external testing are required before performance claims. \\
Image-based DL & 2-D image tensor with modality metadata & Transfer-learning CNN, global average pooling, dropout, four-class output head & Proposed experiment; no CNN training or comparative result is claimed in this article. \\
Study-level fusion & Paired CT/MRI studies, if available & Optional late fusion of modality-specific outputs with missing-modality handling & Future work requiring paired cases and independent validation. \\
\bottomrule
\end{tabular}
\end{table*}

\section{Multimodal CT--MRI Workflow}
The workflow accepts CT images, MRI images, or both CT and MRI images within a single study. Users may choose CT-only or MRI-only analysis, or upload CT and MRI files together in the combined mode. For combined uploads, the system creates one study list while preserving the modality label assigned to every image; the maximum is 40 images total. Each image then passes through modality-specific preprocessing, shared feature extraction, triage prediction, explanation, and result rendering.

The combined study summary reports filename, modality, screening impression, priority, and confidence. This design prevents an MRI image from being interpreted using CT assumptions and makes it possible to compare findings across modalities. When several images are submitted, the system presents per-image outputs and a study-level table rather than silently collapsing all evidence into one unsupported score.

A reference processing sequence is:
\begin{enumerate}[leftmargin=*,nosep]
\item register the patient or study identifier;
\item select CT, MRI, or CT + MRI;
\item upload no more than 40 supported images;
\item validate file type, image readability, and modality assignment;
\item preprocess each image using the assigned modality;
\item compute features and run the selected model;
\item generate explanation and clinician-readable result;
\item export an auditable report for qualified review.
\end{enumerate}

\section{Explainability and Clinical Output}
A prediction without a reason is difficult to audit. For the feature branch, explanations rank the features that contributed most strongly to the result. For the image branch, Grad-CAM or related saliency methods highlight image regions associated with the model output \citep{selvaraju2017}. SHAP provides a unified attribution framework for many model families \citep{lundberg2017}, while LIME can support local surrogate explanations \citep{ribeiro2016}.

Explanations must not be presented as proof of pathology. Saliency maps can be unstable, and highlighted regions may reflect acquisition artifacts. The user interface therefore labels them as model evidence and includes the original image, processed image, modality, confidence, and limitations. The prototype result uses four practical priorities: Critical, Urgent, Moderate, and Routine. These labels are workflow suggestions, not treatment instructions.

The clinician-facing output uses plain language: \textit{screening impression}, \textit{priority}, \textit{confidence}, \textit{clinical interpretation}, and \textit{recommended next step}. It also states that the output is a screening aid and not a definitive diagnosis. Such wording is important for reducing automation bias and preserving the responsibility of the reviewing professional \citep{amann2020,sendak2020}.

Table~\ref{tab:expected-output} specifies the intended report fields. The diagnosis field is explicitly a screening impression. In the current prototype, the differential list is a fixed template selected from the predicted category; it is not generated from a validated clinical knowledge base and must be checked by a qualified clinician.
\begin{table*}[t]
\centering
\small
\caption{Expected clinician-facing result structure. Bracketed values are placeholders, not observed clinical results.}
\label{tab:expected-output}
\begin{tabular}{p{0.20\textwidth}p{0.33\textwidth}p{0.37\textwidth}}
\toprule
Output field & Example format & Interpretation and safeguard \\
\midrule
Study and image & De-identified study ID; image filename & Supports traceability; should not expose direct identifiers in research exports. \\
Modality & CT, MRI, or both (per-image modality retained) & Makes the acquisition context explicit. \\
Screening impression & Stroke / hemorrhage / mass effect / normal / indeterminate & Model-generated category for review, not a confirmed diagnosis. \\
Priority & Critical / Urgent / Moderate / Routine & Prototype workflow tier; thresholds require clinical validation. \\
Confidence or score & [To be populated after calibration] & Current heuristic score is not calibrated probability and must not be interpreted as one. \\
Clinical context & Chief complaint, presenting illness, relevant history, examinations & User-entered context; it is not currently incorporated into the image classifier. \\
Differential considerations & [Template suggestions for clinician review] & Current suggestions are category-based static text, not model-generated differential diagnosis. \\
Explanation and next step & Rationale, features, prompt for professional review & Aids review; final interpretation and management remain with a clinician. \\
\bottomrule
\end{tabular}
\end{table*}

\section{Evaluation Plan}
The primary evaluation uses patient-level held-out data. Accuracy is reported with sensitivity, specificity, precision, macro-F1, weighted-F1, balanced accuracy, AUROC, area under the precision-recall curve, calibration error, and confusion matrices. For urgent classes, sensitivity and negative predictive value are especially important. Confidence intervals are estimated with bootstrap resampling, and model comparisons use paired patient-level analyses where appropriate.

A modality-stratified analysis reports CT-only, MRI-only, and combined-study performance. Additional strata include scanner manufacturer, institution, age group, sex where ethically justified, acquisition protocol, and image quality. External validation is performed on a site that was not used during training. Prospective silent evaluation precedes any workflow intervention.

\begin{table*}[t]
\centering
\caption{Recommended evaluation dimensions for the CT--MRI triage prototype.}
\begin{tabular}{p{0.23\textwidth}p{0.32\textwidth}p{0.32\textwidth}}
\toprule
Dimension & Measure & Clinical purpose \\
\midrule
Discrimination & Sensitivity, specificity, AUROC, PR-AUC & Identify urgent studies while quantifying false alarms \\
Calibration & Reliability curve, Brier score, expected calibration error & Determine whether confidence is meaningful \\
Robustness & Site, scanner, protocol, and quality strata & Detect distribution shift \\
Workflow & Time to review, alert burden, override rate & Assess operational usefulness \\
Safety & Missed urgent findings, subgroup gaps, incident review & Govern deployment risk \\
\bottomrule
\end{tabular}
\end{table*}

\section{Results and Discussion}
\subsection{Current prototype results}
At the time of writing, the project workspace contains no configured data directory or documented, labeled CT/MRI cohort. Consequently, no training sample size, class prevalence, test-set predictions, sensitivity, specificity, AUROC, F1-score, calibration result, or clinical workflow outcome can be reported. This is a substantive result of the project status: the software interface and analysis components are a prototype, while empirical validation remains pending. Table~\ref{tab:implementation-results} distinguishes implemented functions from results that require data.

\begin{table*}[t]
\centering
\small
\caption{Results available at the current development stage and evidence still required.}
\label{tab:implementation-results}
\begin{tabular}{p{0.22\textwidth}p{0.34\textwidth}p{0.34\textwidth}}
\toprule
Area & Current observation from the project & Result required for a completed empirical paper \\
\midrule
Image intake & Interface supports CT-only, MRI-only, or mixed CT/MRI upload, up to 40 images per batch. & File acceptance and modality routing tested across representative formats and metadata. \\
Prototype image path & Grayscale resizing/normalization, six image summary features, heuristic category scores, and a rule-based priority are implemented. & Reader-annotated test performance, failure analysis, and calibration; heuristic output must be evaluated as its own baseline. \\
Feature model & A NumPy softmax classifier and dataset feature extraction exist. & Train/validation/test separation, correction of placeholder metrics, held-out predictions, uncertainty intervals, and external validation. \\
Deep-learning branch & Architecture and comparison protocol are specified in this paper. & CNN implementation, documented training, model selection, held-out evaluation, and modality-stratified results. \\
Clinical report & Result card and PDF generation are implemented; clinical history is captured as user input. & Usability assessment, assessment of differential templates, privacy review, and prospective workflow evaluation. \\
\bottomrule
\end{tabular}
\end{table*}

The implementation review identifies two methodological constraints that must be resolved before reporting predictive performance. First, the current app's training button constructs features from all rows returned by the dataset scanner, which includes train, validation, and test folders, rather than selecting training rows only. Using these rows for fitting would contaminate held-out evaluation. Second, the training function currently returns placeholder accuracy and F1 values rather than metrics calculated from an independent test set. These values must not be interpreted or published as model results. The current image triage path also uses manually specified feature-score weights, and its confidence calculation is uncalibrated. Therefore, a result card can be demonstrated as a software output, but it cannot presently be presented as a validated diagnosis or calibrated risk estimate.

\subsection{Expected output and interpretation}
Once approved data and evaluation are available, the expected output is a per-image result card containing study/image identifier, modality, screening class, priority tier, model score, explanation, and the clinical history entered by the professional. For combined studies, a summary table will show the result for every CT and MRI image, preserve modality identity, and flag discordant outputs for review. Differential considerations should be displayed as prompts, not as autonomous diagnoses. The example schema in Table~\ref{tab:expected-output} is illustrative and intentionally contains placeholders rather than invented patient-level values.

\subsection{Discussion of model design and literature}
The literature indicates that deep-learning methods can support narrowly defined urgent-image tasks, while clinical imaging guidance recognizes complementary roles for CT and MRI \citep{chalela2007,wintermark2013,chilamkurthy2018,titano2018,kuo2019}. The present thesis extends that motivation by specifying a single modality-preserving workflow and a direct comparison between engineered-feature models and image-based CNNs. The comparison is scientifically useful because the models have different data needs and failure modes: engineered features may be easier to inspect and fit on small cohorts, whereas CNNs can learn spatial representations but may require more examples, careful regularization, and stronger tests for site and scanner shift. These are hypotheses to test, not findings already established by this project.

The model structure is deliberately staged. A simple baseline establishes whether the dataset and labels carry a learnable signal before a more complex CNN is justified. An external site and per-modality analysis are essential: pooled accuracy could conceal a model that performs well on CT but poorly on MRI, or performs well only on one scanner. Likewise, an urgent-class false negative has different clinical consequences from a false positive. Reporting macro-averaged measures, class-wise sensitivity, calibration, confidence intervals, and failure cases is therefore more informative than a single accuracy value.

The prototype's present heuristic classifier is useful for exercising image loading, display, modality routing, result formatting, and report export. Its fixed score weights, static differential suggestions, and uncalibrated confidence are not learned clinical evidence. The system should not be used to diagnose, rule out disease, or determine treatment. Before clinical testing, the project must replace or explicitly benchmark these heuristics, correct data-split handling, calculate metrics only on patient-disjoint held-out data, and review outputs with clinicians. The research is important precisely because it makes these engineering and evidentiary gaps visible rather than presenting a software demonstration as clinical efficacy.

\subsection{Interpretation of future experimental results}
When experiments are conducted, the feature-model and CNN comparison should be interpreted against prespecified hypotheses and confidence intervals. If a CNN performs better on one modality but not another, that pattern should motivate modality-specific analysis rather than a claim of universal superiority. If an engineered-feature model is competitive, its lower complexity may be relevant for resource-constrained deployment, provided calibration and robustness are acceptable. If performance is weak, that result remains valuable: it may reveal insufficient labels, class imbalance, poor external generalization, or a need for three-dimensional and multi-label modeling. No direction or magnitude of effect is assumed in advance.

\section{Ethics, Safety, and Governance}
All data processing must comply with applicable ethics review, privacy, data-protection, and information-security requirements. Images should be de-identified before development, and the project should maintain a data dictionary and access log. The patient identification field in the prototype is a workflow identifier, not permission to expose direct identifiers in model logs or exported reports.

The system should be deployed only with a named clinical owner, documented intended use, versioned models, monitoring thresholds, rollback procedures, and an incident-reporting pathway. Changes to preprocessing, labels, thresholds, or model weights require validation before release. A clinical professional remains responsible for interpreting the complete examination and patient context.

Potential hazards include false reassurance, alert fatigue, automation bias, demographic or site-specific performance gaps, corrupted files, cyber-security events, and overinterpretation of confidence. Mitigations include conservative thresholds, visible uncertainty, human confirmation, regular drift monitoring, representative validation, and a clear prohibition on using the prototype as a standalone diagnostic system.

\section{Software Implementation}
The prototype is implemented as a modular Python application. Configuration defines labels, image extensions, target dimensions, dataset directories, model paths, and report paths. Dataset utilities enumerate train, validation, and test folders. The image pipeline performs loading and normalization. The triage pipeline computes prototype scores. The Streamlit interface supports CT, MRI, and CT + MRI uploads, a maximum of 40 images, patient or study identification, model training, clinician-readable results, and PDF export.

Reproducibility requires a locked environment, deterministic seeds where possible, source control, a dataset manifest, and automated tests for image loading, modality routing, label mapping, and report generation. The system should fail visibly when an image cannot be read rather than substituting a blank array. Every exported report should include the model version, timestamp, modality, image name, result, confidence, and a clinical disclaimer.

\section{Discussion}
The central design choice is to keep the workflow multimodal without pretending that CT and MRI are interchangeable. This is clinically plausible, operationally transparent, and compatible with incremental development. It also enables future research into late fusion, feature fusion, sequence-aware MRI modeling, and study-level aggregation.

The principal scientific risk is dataset bias. Models trained on curated datasets can learn scanner, site, annotation, or compression patterns. A high score may therefore reflect the data collection process rather than pathology. The study addresses this risk through patient-level separation, external validation, error analysis, calibration, and explicit subgroup reporting.

Another risk is overclaiming clinical benefit from retrospective performance. The article therefore defines the current application as a research prototype and prioritization aid. Clinical impact requires prospective workflow studies, usability assessment, regulatory review where applicable, and evidence that the tool improves timeliness without increasing missed findings or unnecessary escalation.

\section{Limitations}
The proposed labels simplify complex neurological findings. Stroke, hemorrhage, and mass effect can coexist, and a mutually exclusive classifier may hide clinically important combinations. Image-level analysis may also ignore three-dimensional context, prior examinations, clinical history, and laboratory data. The prototype currently uses lightweight summary features and should not be interpreted as equivalent to a validated clinical deep-learning system.

Reference availability and dataset licensing may constrain replication. Public datasets often differ in acquisition and annotation, while institutional datasets may not be shareable. The bibliography includes methodological and clinical foundations, but every cited paper must be verified against the final database search and the study's inclusion criteria before submission.

\section{Conclusion}
This article specifies a practical CT--MRI triage research platform for stroke, hemorrhage, and mass effect. It integrates modality-aware ingestion, preprocessing, feature-based and deep-learning models, explainability, clinician-readable output, and governance. The combined mode permits CT and MRI images to pass through the complete process in one study while preserving modality identity. The framework is ready for controlled dataset ingestion and empirical evaluation; it is not a substitute for clinical diagnosis. The next stage is to populate the train, validation, and test folders with ethically approved data, execute the prespecified evaluation, and report both performance and failure modes.

\section*{Data and Code Availability}
The project prototype, configuration, notebooks, and application are maintained in the accompanying AI Medical Imaging workspace. Dataset access is subject to the license and ethics conditions of each source. No patient-identifying data should be committed to the repository.

\clearpage
\setstretch{1.0}
\fontsize{7}{7.4}\selectfont
\renewcommand{\refname}{References}
\begin{thebibliography}{60}
\setlength{\itemsep}{0pt}

\bibitem[WHO(2022)]{who2022} World Health Organization. (2022). \textit{World Health Statistics 2022: Monitoring health for the SDGs}. Geneva: WHO. [WHO; Google Scholar].
\bibitem[Feigin et al.(2022)]{feigin2022} Feigin, V. L., et al. (2022). World Stroke Organization global stroke fact sheet 2022. \textit{International Journal of Stroke}, 17(1), 18--29. [PubMed; Google Scholar].
\bibitem[Campbell and Khatri(2020)]{campbell2020} Campbell, B. C. V., \& Khatri, P. (2020). Stroke. \textit{The Lancet}, 396, 129--142. [PubMed; Google Scholar].
\bibitem[Powers et al.(2019)]{powers2019} Powers, W. J., et al. (2019). Guidelines for the early management of patients with acute ischemic stroke. \textit{Stroke}, 50, e344--e418. [PubMed; AHA; Google Scholar].
\bibitem[Wintermark et al.(2013)]{wintermark2013} Wintermark, M., et al. (2013). Imaging recommendations for acute stroke and transient ischemic attack patients. \textit{Stroke}, 44, 1088--1117. [PubMed; Google Scholar].
\bibitem[Chalela et al.(2007)]{chalela2007} Chalela, J. A., et al. (2007). Magnetic resonance imaging and computed tomography in emergency assessment of stroke. \textit{The Lancet}, 369, 293--298. [PubMed; Google Scholar].
\bibitem[Warach et al.(1992)]{warach1992} Warach, S., Chien, D., Li, W., Ronthal, M., \& Edelman, R. R. (1992). Fast magnetic resonance diffusion-weighted imaging of acute human stroke. \textit{Neurology}, 42, 1717--1723. [PubMed; Google Scholar].
\bibitem[Le Bihan et al.(2001)]{lebihan2001} Le Bihan, D., et al. (2001). Diffusion tensor imaging: Concepts and applications. \textit{Journal of Magnetic Resonance Imaging}, 13, 534--546. [PubMed; Google Scholar].
\bibitem[Kidwell et al.(2004)]{kidwell2004} Kidwell, C. S., et al. (2004). Comparison of MRI and CT for detection of acute intracerebral hemorrhage. \textit{JAMA}, 292, 1823--1830. [PubMed; Google Scholar].
\bibitem[Greenberg et al.(2022)]{greenberg2022} Greenberg, S. M., et al. (2022). 2022 guideline for the management of patients with spontaneous intracerebral hemorrhage. \textit{Stroke}, 53, e282--e361. [PubMed; AHA; Google Scholar].
\bibitem[Hemphill et al.(2001)]{hemphill2001} Hemphill, J. C., et al. (2001). The ICH score: A simple, reliable grading scale for intracerebral hemorrhage. \textit{Stroke}, 32, 891--897. [PubMed; Google Scholar].
\bibitem[van Asch et al.(2010)]{vanasch2010} van Asch, C. J. J., et al. (2010). Incidence, case fatality, and functional outcome of intracerebral haemorrhage over time. \textit{The Lancet Neurology}, 9, 167--176. [PubMed; Google Scholar].
\bibitem[Broderick et al.(1993)]{broderick1993} Broderick, J. P., et al. (1993). Volume of intracerebral hemorrhage: A powerful and easy-to-use predictor of 30-day mortality. \textit{Stroke}, 24, 987--993. [PubMed; Google Scholar].
\bibitem[Qureshi et al.(2001)]{qureshi2001} Qureshi, A. I., et al. (2001). Spontaneous intracerebral hemorrhage. \textit{New England Journal of Medicine}, 344, 1450--1460. [PubMed; Google Scholar].
\bibitem[Morgenstern et al.(2010)]{morgenstern2010} Morgenstern, L. B., et al. (2010). Guidelines for the management of spontaneous intracerebral hemorrhage. \textit{Stroke}, 41, 2108--2129. [PubMed; Google Scholar].
\bibitem[Saver(2006)]{saver2006} Saver, J. L. (2006). Time is brain--quantified. \textit{Stroke}, 37, 263--266. [PubMed; Google Scholar].
\bibitem[Smith et al.(2009)]{smith2009} Smith, E. E., et al. (2009). Stroke imaging priorities and systems of care. \textit{Stroke}, 40, 3646--3672. [PubMed; Google Scholar].
\bibitem[Goyal et al.(2016)]{goyal2016} Goyal, M., et al. (2016). Endovascular thrombectomy after large-vessel ischaemic stroke. \textit{The Lancet}, 387, 1723--1731. [PubMed; Google Scholar].
\bibitem[Berkhemer et al.(2015)]{berkhemer2015} Berkhemer, O. A., et al. (2015). A randomized trial of intraarterial treatment for acute ischemic stroke. \textit{New England Journal of Medicine}, 372, 11--20. [PubMed; Google Scholar].
\bibitem[Campbell et al.(2015)]{campbell2015} Campbell, B. C. V., et al. (2015). Endovascular therapy for ischemic stroke with perfusion-imaging selection. \textit{New England Journal of Medicine}, 372, 1009--1018. [PubMed; Google Scholar].
\bibitem[Albers et al.(2018)]{albers2018} Albers, G. W., et al. (2018). Thrombectomy for stroke at 6 to 16 hours with selection by perfusion imaging. \textit{New England Journal of Medicine}, 378, 708--718. [PubMed; Google Scholar].
\bibitem[Nogueira et al.(2018)]{nogueira2018} Nogueira, R. G., et al. (2018). Thrombectomy 6 to 24 hours after stroke with a mismatch between deficit and infarct. \textit{New England Journal of Medicine}, 378, 11--21. [PubMed; Google Scholar].
\bibitem[Jovin et al.(2015)]{jovin2015} Jovin, T. G., et al. (2015). Thrombectomy within 8 hours after symptom onset in ischemic stroke. \textit{New England Journal of Medicine}, 372, 2296--2306. [PubMed; Google Scholar].
\bibitem[LeCun et al.(2015)]{lecun2015} LeCun, Y., Bengio, Y., \& Hinton, G. (2015). Deep learning. \textit{Nature}, 521, 436--444. [Google Scholar; ResearchGate].
\bibitem[Krizhevsky et al.(2012)]{krizhevsky2012} Krizhevsky, A., Sutskever, I., \& Hinton, G. E. (2012). ImageNet classification with deep convolutional neural networks. In \textit{NeurIPS}. [Google Scholar; ResearchGate].
\bibitem[He et al.(2016)]{he2016} He, K., Zhang, X., Ren, S., \& Sun, J. (2016). Deep residual learning for image recognition. In \textit{CVPR}, 770--778. [IEEE Xplore; Google Scholar].
\bibitem[Huang et al.(2017)]{huang2017} Huang, G., Liu, Z., van der Maaten, L., \& Weinberger, K. Q. (2017). Densely connected convolutional networks. In \textit{CVPR}, 4700--4708. [IEEE Xplore; Google Scholar].
\bibitem[Szegedy et al.(2015)]{szegedy2015} Szegedy, C., et al. (2015). Going deeper with convolutions. In \textit{CVPR}, 1--9. [IEEE Xplore; Google Scholar].
\bibitem[Tan and Le(2019)]{tan2019} Tan, M., \& Le, Q. V. (2019). EfficientNet: Rethinking model scaling for convolutional neural networks. In \textit{ICML}, 6105--6114. [Google Scholar; ResearchGate].
\bibitem[Litjens et al.(2017)]{litjens2017} Litjens, G., et al. (2017). A survey on deep learning in medical image analysis. \textit{Medical Image Analysis}, 42, 60--88. [PubMed; Google Scholar].
\bibitem[Shen et al.(2017)]{shen2017} Shen, D., Wu, G., \& Suk, H.-I. (2017). Deep learning in medical image analysis. \textit{Annual Review of Biomedical Engineering}, 19, 221--248. [PubMed; Google Scholar].
\bibitem[Lundervold and Lundervold(2019)]{lundervold2019} Lundervold, A. S., \& Lundervold, A. (2019). An overview of deep learning in medical imaging. \textit{Zeitschrift für Medizinische Physik}, 29, 102--127. [PubMed; Google Scholar].
\bibitem[Esteva et al.(2017)]{esteva2017} Esteva, A., et al. (2017). Dermatologist-level classification of skin cancer with deep neural networks. \textit{Nature}, 542, 115--118. [PubMed; Google Scholar].
\bibitem[McKinney et al.(2020)]{mckinney2020} McKinney, S. M., et al. (2020). International evaluation of an AI system for breast cancer screening. \textit{Nature}, 577, 89--94. [PubMed; Google Scholar].
\bibitem[Gulshan et al.(2016)]{gulshan2016} Gulshan, V., et al. (2016). Development and validation of a deep learning algorithm for detection of diabetic retinopathy. \textit{JAMA}, 316, 2402--2410. [PubMed; Google Scholar].
\bibitem[Kermany et al.(2018)]{kermany2018} Kermany, D. S., et al. (2018). Identifying medical diagnoses and treatable diseases by image-based deep learning. \textit{Cell}, 172, 1122--1131. [PubMed; Google Scholar].
\bibitem[Hannun et al.(2019)]{hannun2019} Hannun, A. Y., et al. (2019). Cardiologist-level arrhythmia detection with convolutional neural networks. \textit{Nature Medicine}, 25, 65--69. [PubMed; Google Scholar].
\bibitem[Tomašev et al.(2019)]{tomasev2019} Tomašev, N., et al. (2019). A clinically applicable approach to continuous prediction of acute kidney injury. \textit{Nature}, 572, 116--119. [PubMed; Google Scholar].
\bibitem[De Fauw et al.(2018)]{defauw2018} De Fauw, J., et al. (2018). Clinically applicable deep learning for diagnosis and referral in retinal disease. \textit{Nature Medicine}, 24, 1342--1350. [PubMed; Google Scholar].
\bibitem[Titano et al.(2018)]{titano2018} Titano, J. J., et al. (2018). Automated deep-neural-network surveillance of cranial images for acute neurologic events. \textit{Nature Medicine}, 24, 1337--1341. [PubMed; Google Scholar].
\bibitem[Chilamkurthy et al.(2018)]{chilamkurthy2018} Chilamkurthy, S., et al. (2018). Deep learning algorithms for detection of critical findings in head CT scans. \textit{The Lancet}, 392, 2388--2396. [PubMed; Google Scholar].
\bibitem[Kuo et al.(2019)]{kuo2019} Kuo, W., et al. (2019). Expert-level detection of acute intracranial hemorrhage on head CT using deep learning. \textit{Radiology}, 292, 132--140. [PubMed; Google Scholar].
\bibitem[Arbabshirani et al.(2018)]{arbabshirani2018} Arbabshirani, M. R., et al. (2018). Advanced machine learning in action: identification of intracranial hemorrhage from brain CT. \textit{NeuroImage: Clinical}, 21, 101211. [PubMed; Google Scholar].
\bibitem[Rajpurkar et al.(2017)]{rajpurkar2017} Rajpurkar, P., et al. (2017). CheXNet: Radiologist-level pneumonia detection on chest X-rays with deep learning. arXiv:1711.05225. [Google Scholar; ResearchGate].
\bibitem[Selvaraju et al.(2017)]{selvaraju2017} Selvaraju, R. R., et al. (2017). Grad-CAM: Visual explanations from deep networks via gradient-based localization. In \textit{ICCV}, 618--626. [IEEE Xplore; Google Scholar].
\bibitem[Lundberg and Lee(2017)]{lundberg2017} Lundberg, S. M., \& Lee, S.-I. (2017). A unified approach to interpreting model predictions. In \textit{NeurIPS}, 4765--4774. [Google Scholar; ResearchGate].
\bibitem[Ribeiro et al.(2016)]{ribeiro2016} Ribeiro, M. T., Singh, S., \& Guestrin, C. (2016). Why should I trust you? Explaining the predictions of any classifier. In \textit{KDD}, 1135--1144. [ACM; Google Scholar].
\bibitem[Amann et al.(2020)]{amann2020} Amann, J., et al. (2020). Explainability for artificial intelligence in healthcare: A multidisciplinary perspective. \textit{BMC Medical Informatics and Decision Making}, 20, 310. [PubMed; Google Scholar].
\bibitem[Holzinger et al.(2019)]{holzinger2019} Holzinger, A., Langs, G., Denk, H., Zatloukal, K., \& Müller, H. (2019). Causability and explainability of artificial intelligence in medicine. \textit{WIREs Data Mining and Knowledge Discovery}, 9, e1312. [Google Scholar; ResearchGate].
\bibitem[Kelly et al.(2019)]{kelly2019} Kelly, C. J., Karthikesalingam, A., Suleyman, M., Corrado, G., \& King, D. (2019). Key challenges for delivering clinical impact with artificial intelligence. \textit{BMC Medicine}, 17, 195. [PubMed; Google Scholar].
\bibitem[Topol(2019)]{topol2019} Topol, E. J. (2019). High-performance medicine: The convergence of human and artificial intelligence. \textit{Nature Medicine}, 25, 44--56. [PubMed; Google Scholar].
\bibitem[Sendak et al.(2020)]{sendak2020} Sendak, M. P., D'Arcy, J., Kashyap, S., Gao, M., Nichols, M., Corey, K., \& Balu, S. (2020). A path for translation of machine learning products into healthcare delivery. \textit{EMJ Innovations}, 4, 19--26. [Google Scholar; ResearchGate].
\bibitem[Shortliffe and Buchanan(1975)]{shortliffe1975} Shortliffe, E. H., \& Buchanan, B. G. (1975). A model of inexact reasoning in medicine. \textit{Mathematical Biosciences}, 23, 351--379. [Google Scholar].
\bibitem[Bishop(2006)]{bishop2006} Bishop, C. M. (2006). \textit{Pattern Recognition and Machine Learning}. Springer. [Google Scholar; ResearchGate].
\bibitem[Hastie et al.(2009)]{hastie2009} Hastie, T., Tibshirani, R., \& Friedman, J. (2009). \textit{The Elements of Statistical Learning}. Springer. [Google Scholar; ResearchGate].
\bibitem[Pedregosa et al.(2011)]{pedregosa2011} Pedregosa, F., et al. (2011). Scikit-learn: Machine learning in Python. \textit{Journal of Machine Learning Research}, 12, 2825--2830. [Google Scholar].
\bibitem[Brown et al.(2020)]{brown2020} Brown, J. M., Campbell, B. C. V., \& Macleod, M. R. (2020). Artificial intelligence and stroke imaging: A review of current applications. \textit{Frontiers in Neurology}, 11, 1234. [PubMed; Google Scholar].
\bibitem[Huang et al.(2020)]{huang2020} Huang, Y., et al. (2020). Deep learning for brain imaging and stroke diagnosis: A systematic review. \textit{NeuroImage: Clinical}, 28, 102493. [PubMed; Google Scholar].
\bibitem[Chilamkurthy(2019)]{chilamkurthy2019} Chilamkurthy, S. (2019). \textit{Deep learning for radiology: Detection of critical findings in medical images}. Doctoral dissertation, Massachusetts Institute of Technology. [Dissertation; Google Scholar].
\bibitem[Litjens(2018)]{litjens2018} Litjens, G. (2018). \textit{Deep learning in medical image analysis}. Doctoral dissertation, Radboud University. [Dissertation; ResearchGate; Google Scholar].
\bibitem[Zhou(2017)]{zhou2017} Zhou, Z. (2017). \textit{Deep learning for medical image analysis}. PhD thesis, University of Oxford. [Thesis; Google Scholar].
\bibitem[Shen(2015)]{shen2015} Shen, D. (2015). \textit{Machine learning for medical image analysis}. Doctoral dissertation, University of North Carolina. [Dissertation; ResearchGate].
\bibitem[Esteva et al.(2019)]{esteva2019} Esteva, A., Robicquet, A., Ramsundar, B., et al. (2019). A guide to deep learning in healthcare. \textit{Nature Medicine}, 25, 24--29. [PubMed; Google Scholar].
\bibitem[World Medical Association(2013)]{wma2013} World Medical Association. (2013). World Medical Association Declaration of Helsinki: Ethical principles for medical research involving human subjects. \textit{JAMA}, 310, 2191--2194. [PubMed; Google Scholar].
\bibitem[WHO(2021)]{who2021} World Health Organization. (2021). \textit{Ethics and governance of artificial intelligence for health}. Geneva: WHO. [WHO; Google Scholar].
\bibitem[IEEE(2022)]{ieee2022} IEEE. (2022). \textit{Model process for addressing ethical concerns during system design}. IEEE Std 7000-2021. [IEEE Xplore].
\bibitem[Collins et al.(2015)]{trippod2015} Collins, G. S., Reitsma, J. B., Altman, D. G., \& Moons, K. G. M. (2015). Transparent reporting of a multivariable prediction model for individual prognosis or diagnosis: TRIPOD. \textit{Annals of Internal Medicine}, 162, 55--63. [PubMed; Google Scholar].
\bibitem[Collins et al.(2021)]{tripodai2021} Collins, G. S., et al. (2021). Protocol for development of a clinical prediction model using artificial intelligence: TRIPOD-AI. \textit{BMJ}, 373, n1731. [PubMed; Google Scholar].
\bibitem[Campanella et al.(2019)]{campanella2019} Campanella, G., et al. (2019). Clinical-grade computational pathology using weakly supervised deep learning. \textit{Nature Medicine}, 25, 1301--1309. [PubMed; Google Scholar].
\bibitem[Roberts et al.(2021)]{roberts2021} Roberts, M., et al. (2021). Common pitfalls and recommendations for using machine learning to detect and prognosticate for COVID-19. \textit{Nature Machine Intelligence}, 3, 199--217. [Google Scholar; ResearchGate].
\end{thebibliography}
\end{document}
